% Options for packages loaded elsewhere
% Options for packages loaded elsewhere
\PassOptionsToPackage{unicode}{hyperref}
\PassOptionsToPackage{hyphens}{url}
\PassOptionsToPackage{dvipsnames,svgnames,x11names}{xcolor}
%
\documentclass[
  9pt,
  letterpaper,
  DIV=11,
  numbers=noendperiod]{scrartcl}
\usepackage{xcolor}
\usepackage[margin=0.4in]{geometry}
\usepackage{amsmath,amssymb}
\setcounter{secnumdepth}{-\maxdimen} % remove section numbering
\usepackage{iftex}
\ifPDFTeX
  \usepackage[T1]{fontenc}
  \usepackage[utf8]{inputenc}
  \usepackage{textcomp} % provide euro and other symbols
\else % if luatex or xetex
  \usepackage{unicode-math} % this also loads fontspec
  \defaultfontfeatures{Scale=MatchLowercase}
  \defaultfontfeatures[\rmfamily]{Ligatures=TeX,Scale=1}
\fi
\usepackage{lmodern}
\ifPDFTeX\else
  % xetex/luatex font selection
\fi
% Use upquote if available, for straight quotes in verbatim environments
\IfFileExists{upquote.sty}{\usepackage{upquote}}{}
\IfFileExists{microtype.sty}{% use microtype if available
  \usepackage[]{microtype}
  \UseMicrotypeSet[protrusion]{basicmath} % disable protrusion for tt fonts
}{}
\usepackage{setspace}
\makeatletter
\@ifundefined{KOMAClassName}{% if non-KOMA class
  \IfFileExists{parskip.sty}{%
    \usepackage{parskip}
  }{% else
    \setlength{\parindent}{0pt}
    \setlength{\parskip}{6pt plus 2pt minus 1pt}}
}{% if KOMA class
  \KOMAoptions{parskip=half}}
\makeatother
% Make \paragraph and \subparagraph free-standing
\makeatletter
\ifx\paragraph\undefined\else
  \let\oldparagraph\paragraph
  \renewcommand{\paragraph}{
    \@ifstar
      \xxxParagraphStar
      \xxxParagraphNoStar
  }
  \newcommand{\xxxParagraphStar}[1]{\oldparagraph*{#1}\mbox{}}
  \newcommand{\xxxParagraphNoStar}[1]{\oldparagraph{#1}\mbox{}}
\fi
\ifx\subparagraph\undefined\else
  \let\oldsubparagraph\subparagraph
  \renewcommand{\subparagraph}{
    \@ifstar
      \xxxSubParagraphStar
      \xxxSubParagraphNoStar
  }
  \newcommand{\xxxSubParagraphStar}[1]{\oldsubparagraph*{#1}\mbox{}}
  \newcommand{\xxxSubParagraphNoStar}[1]{\oldsubparagraph{#1}\mbox{}}
\fi
\makeatother

\usepackage{color}
\usepackage{fancyvrb}
\newcommand{\VerbBar}{|}
\newcommand{\VERB}{\Verb[commandchars=\\\{\}]}
\DefineVerbatimEnvironment{Highlighting}{Verbatim}{commandchars=\\\{\}}
% Add ',fontsize=\small' for more characters per line
\usepackage{framed}
\definecolor{shadecolor}{RGB}{241,243,245}
\newenvironment{Shaded}{\begin{snugshade}}{\end{snugshade}}
\newcommand{\AlertTok}[1]{\textcolor[rgb]{0.68,0.00,0.00}{#1}}
\newcommand{\AnnotationTok}[1]{\textcolor[rgb]{0.37,0.37,0.37}{#1}}
\newcommand{\AttributeTok}[1]{\textcolor[rgb]{0.40,0.45,0.13}{#1}}
\newcommand{\BaseNTok}[1]{\textcolor[rgb]{0.68,0.00,0.00}{#1}}
\newcommand{\BuiltInTok}[1]{\textcolor[rgb]{0.00,0.23,0.31}{#1}}
\newcommand{\CharTok}[1]{\textcolor[rgb]{0.13,0.47,0.30}{#1}}
\newcommand{\CommentTok}[1]{\textcolor[rgb]{0.37,0.37,0.37}{#1}}
\newcommand{\CommentVarTok}[1]{\textcolor[rgb]{0.37,0.37,0.37}{\textit{#1}}}
\newcommand{\ConstantTok}[1]{\textcolor[rgb]{0.56,0.35,0.01}{#1}}
\newcommand{\ControlFlowTok}[1]{\textcolor[rgb]{0.00,0.23,0.31}{\textbf{#1}}}
\newcommand{\DataTypeTok}[1]{\textcolor[rgb]{0.68,0.00,0.00}{#1}}
\newcommand{\DecValTok}[1]{\textcolor[rgb]{0.68,0.00,0.00}{#1}}
\newcommand{\DocumentationTok}[1]{\textcolor[rgb]{0.37,0.37,0.37}{\textit{#1}}}
\newcommand{\ErrorTok}[1]{\textcolor[rgb]{0.68,0.00,0.00}{#1}}
\newcommand{\ExtensionTok}[1]{\textcolor[rgb]{0.00,0.23,0.31}{#1}}
\newcommand{\FloatTok}[1]{\textcolor[rgb]{0.68,0.00,0.00}{#1}}
\newcommand{\FunctionTok}[1]{\textcolor[rgb]{0.28,0.35,0.67}{#1}}
\newcommand{\ImportTok}[1]{\textcolor[rgb]{0.00,0.46,0.62}{#1}}
\newcommand{\InformationTok}[1]{\textcolor[rgb]{0.37,0.37,0.37}{#1}}
\newcommand{\KeywordTok}[1]{\textcolor[rgb]{0.00,0.23,0.31}{\textbf{#1}}}
\newcommand{\NormalTok}[1]{\textcolor[rgb]{0.00,0.23,0.31}{#1}}
\newcommand{\OperatorTok}[1]{\textcolor[rgb]{0.37,0.37,0.37}{#1}}
\newcommand{\OtherTok}[1]{\textcolor[rgb]{0.00,0.23,0.31}{#1}}
\newcommand{\PreprocessorTok}[1]{\textcolor[rgb]{0.68,0.00,0.00}{#1}}
\newcommand{\RegionMarkerTok}[1]{\textcolor[rgb]{0.00,0.23,0.31}{#1}}
\newcommand{\SpecialCharTok}[1]{\textcolor[rgb]{0.37,0.37,0.37}{#1}}
\newcommand{\SpecialStringTok}[1]{\textcolor[rgb]{0.13,0.47,0.30}{#1}}
\newcommand{\StringTok}[1]{\textcolor[rgb]{0.13,0.47,0.30}{#1}}
\newcommand{\VariableTok}[1]{\textcolor[rgb]{0.07,0.07,0.07}{#1}}
\newcommand{\VerbatimStringTok}[1]{\textcolor[rgb]{0.13,0.47,0.30}{#1}}
\newcommand{\WarningTok}[1]{\textcolor[rgb]{0.37,0.37,0.37}{\textit{#1}}}

\usepackage{longtable,booktabs,array}
\newcounter{none} % for unnumbered tables
\usepackage{calc} % for calculating minipage widths
% Correct order of tables after \paragraph or \subparagraph
\usepackage{etoolbox}
\makeatletter
\patchcmd\longtable{\par}{\if@noskipsec\mbox{}\fi\par}{}{}
\makeatother
% Allow footnotes in longtable head/foot
\IfFileExists{footnotehyper.sty}{\usepackage{footnotehyper}}{\usepackage{footnote}}
\makesavenoteenv{longtable}
\usepackage{graphicx}
\makeatletter
\newsavebox\pandoc@box
\newcommand*\pandocbounded[1]{% scales image to fit in text height/width
  \sbox\pandoc@box{#1}%
  \Gscale@div\@tempa{\textheight}{\dimexpr\ht\pandoc@box+\dp\pandoc@box\relax}%
  \Gscale@div\@tempb{\linewidth}{\wd\pandoc@box}%
  \ifdim\@tempb\p@<\@tempa\p@\let\@tempa\@tempb\fi% select the smaller of both
  \ifdim\@tempa\p@<\p@\scalebox{\@tempa}{\usebox\pandoc@box}%
  \else\usebox{\pandoc@box}%
  \fi%
}
% Set default figure placement to htbp
\def\fps@figure{htbp}
\makeatother





\setlength{\emergencystretch}{3em} % prevent overfull lines

\providecommand{\tightlist}{%
  \setlength{\itemsep}{0pt}\setlength{\parskip}{0pt}}



 


\KOMAoption{captions}{tableheading}
\makeatletter
\@ifpackageloaded{caption}{}{\usepackage{caption}}
\AtBeginDocument{%
\ifdefined\contentsname
  \renewcommand*\contentsname{Table of contents}
\else
  \newcommand\contentsname{Table of contents}
\fi
\ifdefined\listfigurename
  \renewcommand*\listfigurename{List of Figures}
\else
  \newcommand\listfigurename{List of Figures}
\fi
\ifdefined\listtablename
  \renewcommand*\listtablename{List of Tables}
\else
  \newcommand\listtablename{List of Tables}
\fi
\ifdefined\figurename
  \renewcommand*\figurename{Figure}
\else
  \newcommand\figurename{Figure}
\fi
\ifdefined\tablename
  \renewcommand*\tablename{Table}
\else
  \newcommand\tablename{Table}
\fi
}
\@ifpackageloaded{float}{}{\usepackage{float}}
\floatstyle{ruled}
\@ifundefined{c@chapter}{\newfloat{codelisting}{h}{lop}}{\newfloat{codelisting}{h}{lop}[chapter]}
\floatname{codelisting}{Listing}
\newcommand*\listoflistings{\listof{codelisting}{List of Listings}}
\makeatother
\makeatletter
\makeatother
\makeatletter
\@ifpackageloaded{caption}{}{\usepackage{caption}}
\@ifpackageloaded{subcaption}{}{\usepackage{subcaption}}
\makeatother
\usepackage{bookmark}
\IfFileExists{xurl.sty}{\usepackage{xurl}}{} % add URL line breaks if available
\urlstyle{same}
% fallback for those not using the hyperref driver hyperxmp:
\makeatletter
\@ifundefined{xmpquote}{\newcommand{\xmpquote}[1]{#1}}{}
\makeatother
\hypersetup{
  pdftitle={Data Story Simulation: Do Protest Waves Shift What We Read on Wikipedia?},
  pdfauthor={Brie Beaudry},
  colorlinks=true,
  linkcolor={blue},
  filecolor={Maroon},
  citecolor={Blue},
  urlcolor={Blue},
  pdfcreator={LaTeX via pandoc}}


\title{Data Story Simulation: Do Protest Waves Shift What We Read on
Wikipedia?}
\author{Brie Beaudry}
\date{}
\begin{document}
\maketitle


\setstretch{0.95}
\subsection{Research Question}\label{research-question}

\textbf{How do the Wikipedia pages people click on from the ``Climate
change'' article change during months with major U.S. climate protests?}

Protests bring public attention to particular issues, but it is less
clear whether this attention changes how people seek out information.
This project explores the relationship between collective action in
public spaces and how people navigate environmental knowledge online.

\subsection{Data Sources}\label{data-sources}

\begin{itemize}
\tightlist
\item
  \textbf{Protests:} The
  \href{https://dataverse.harvard.edu/dataset.xhtml?persistentId=doi:10.7910/DVN/6OPP7H}{Crowd
  Counting Consortium (CCC)} collects data on protest events in the
  United States. I will use monthly counts of climate and environmental
  protests, defining a \textbf{protest wave} as a month in the top 10\%
  of counts.
\item
  \textbf{Wikipedia clicks:} The
  \href{https://dumps.wikimedia.org/other/clickstream/}{Wikimedia
  Clickstream dataset} records clicks between Wikipedia pages each
  month. I will focus on links from the ``Climate change'' article to
  see which pages people visit next.
\end{itemize}

\subsection{Simulated Data and
Analysis}\label{simulated-data-and-analysis}

\subsubsection{1: Simulate the
Clickstream}\label{simulate-the-clickstream}

The fake data has 30 destination pages (Page 1 is the most popular, and
popularity falls off quickly) and 8 pairs of months. Each pair has a
``before'' and an ``after'' month. Four pairs are wave pairs, where the
``after'' month is a protest wave. Four are ordinary pairs.

The planted effect is in wave pairs only; Pages 11 to 14 get 75\% more
clicks in the ``after'' month. That makes 30 pages × 8 pairs × 2 months
= 480 rows.

\paragraph{1.1. Setup}\label{setup}

\begin{Shaded}
\begin{Highlighting}[]
\FunctionTok{library}\NormalTok{(tidyverse)}
\FunctionTok{set.seed}\NormalTok{(}\DecValTok{853}\NormalTok{)}
\end{Highlighting}
\end{Shaded}

\paragraph{1.2 Simulating Clickstream}\label{simulating-clickstream}

\begin{Shaded}
\begin{Highlighting}[]
\CommentTok{\#| label: simulate\_data}

\CommentTok{\# Each page gets a "normal" monthly click level}
\NormalTok{pages }\OtherTok{\textless{}{-}} 
  \FunctionTok{tibble}\NormalTok{(}
    \AttributeTok{page =} \FunctionTok{paste}\NormalTok{(}\StringTok{"Page"}\NormalTok{, }\DecValTok{1}\SpecialCharTok{:}\DecValTok{30}\NormalTok{),}
    \AttributeTok{base =} \FunctionTok{round}\NormalTok{(}\DecValTok{20000} \SpecialCharTok{/}\NormalTok{ (}\DecValTok{1}\SpecialCharTok{:}\DecValTok{30}\NormalTok{) }\SpecialCharTok{*} \FunctionTok{runif}\NormalTok{(}\DecValTok{30}\NormalTok{, .}\DecValTok{7}\NormalTok{, }\FloatTok{1.3}\NormalTok{))}
\NormalTok{  )}

\CommentTok{\# Every combination of page x pair x month (30 x 8 x 2 = 480 rows)}
\NormalTok{sim }\OtherTok{\textless{}{-}} \FunctionTok{expand\_grid}\NormalTok{(}
\NormalTok{  pages,}
  \AttributeTok{pair =} \DecValTok{1}\SpecialCharTok{:}\DecValTok{8}\NormalTok{,}
  \AttributeTok{month\_type =} \FunctionTok{c}\NormalTok{(}\StringTok{"before"}\NormalTok{, }\StringTok{"after"}\NormalTok{)}
\NormalTok{) }\SpecialCharTok{|\textgreater{}}
  \FunctionTok{mutate}\NormalTok{(}
    \CommentTok{\# Pairs 1{-}4 are protest waves; pairs 5{-}8 are ordinary}
    \AttributeTok{pair\_type =} \FunctionTok{if\_else}\NormalTok{(pair }\SpecialCharTok{\textless{}=} \DecValTok{4}\NormalTok{, }\StringTok{"wave"}\NormalTok{, }\StringTok{"ordinary"}\NormalTok{),}
    
    \CommentTok{\# Planted effect: +75\% clicks, only for Pages 11{-}14,}
    \CommentTok{\# only in the "after" month, only in wave pairs}
    \AttributeTok{growth =} \FunctionTok{if\_else}\NormalTok{(}
\NormalTok{      pair\_type }\SpecialCharTok{==} \StringTok{"wave"} \SpecialCharTok{\&}\NormalTok{ month\_type }\SpecialCharTok{==} \StringTok{"after"} \SpecialCharTok{\&}
\NormalTok{        page }\SpecialCharTok{\%in\%} \FunctionTok{paste}\NormalTok{(}\StringTok{"Page"}\NormalTok{, }\DecValTok{11}\SpecialCharTok{:}\DecValTok{14}\NormalTok{),}
      \FloatTok{1.75}\NormalTok{, }\DecValTok{1}
\NormalTok{    ),}
    
    \CommentTok{\# Clicks = normal level x growth}
    \AttributeTok{clicks =} \FunctionTok{round}\NormalTok{(base }\SpecialCharTok{*}\NormalTok{ growth)}
\NormalTok{  ) }\SpecialCharTok{|\textgreater{}}
  \FunctionTok{select}\NormalTok{(pair, pair\_type, month\_type, page, clicks)}

\FunctionTok{stopifnot}\NormalTok{(}\FunctionTok{nrow}\NormalTok{(sim) }\SpecialCharTok{==} \DecValTok{480}\NormalTok{)  }\CommentTok{\# safety check}

\FunctionTok{head}\NormalTok{(sim)}
\end{Highlighting}
\end{Shaded}

\begin{verbatim}
# A tibble: 6 x 5
   pair pair_type month_type page   clicks
  <int> <chr>     <chr>      <chr>   <dbl>
1     1 wave      before     Page 1  18314
2     1 wave      after      Page 1  18314
3     2 wave      before     Page 1  18314
4     2 wave      after      Page 1  18314
5     3 wave      before     Page 1  18314
6     3 wave      after      Page 1  18314
\end{verbatim}

\subsection{Step 2: Compare the Top 10 Destination
Pages}\label{step-2-compare-the-top-10-destination-pages}

For each pair and month, I keep the 10 most-clicked pages. Then, for
each pair, I count how many of the ``after'' top 10 were also in the
``before'' top 10. This is the overlap, from 0 to 10. An overlap of 10
means the list did not change. A lower overlap means new pages pushed
into the top 10. Finally, I take the median for wave pairs and for
ordinary pairs.

\begin{Shaded}
\begin{Highlighting}[]
\CommentTok{\# The 10 most{-}clicked pages in each pair and month}
\NormalTok{top10 }\OtherTok{\textless{}{-}}\NormalTok{ sim }\SpecialCharTok{|\textgreater{}}
  \FunctionTok{slice\_max}\NormalTok{(}
\NormalTok{    clicks,}
    \AttributeTok{n =} \DecValTok{10}\NormalTok{,}
    \AttributeTok{by =} \FunctionTok{c}\NormalTok{(pair, month\_type),}
    \AttributeTok{with\_ties =} \ConstantTok{FALSE}
\NormalTok{  )}

\CommentTok{\# For each pair: how many "after" top{-}10 pages were also in the "before" top 10?}
\NormalTok{overlap }\OtherTok{\textless{}{-}}\NormalTok{ top10 }\SpecialCharTok{|\textgreater{}}
  \FunctionTok{group\_by}\NormalTok{(pair, pair\_type) }\SpecialCharTok{|\textgreater{}}
  \FunctionTok{summarise}\NormalTok{(}
    \AttributeTok{shared =} \FunctionTok{sum}\NormalTok{(}
\NormalTok{      page[month\_type }\SpecialCharTok{==} \StringTok{"after"}\NormalTok{] }\SpecialCharTok{\%in\%}
\NormalTok{        page[month\_type }\SpecialCharTok{==} \StringTok{"before"}\NormalTok{]}
\NormalTok{    ),}
    \AttributeTok{.groups =} \StringTok{"drop"}
\NormalTok{  )}

\CommentTok{\# Typical (median) overlap for wave vs. ordinary pairs}
\NormalTok{results }\OtherTok{\textless{}{-}}\NormalTok{ overlap }\SpecialCharTok{|\textgreater{}}
  \FunctionTok{group\_by}\NormalTok{(pair\_type) }\SpecialCharTok{|\textgreater{}}
  \FunctionTok{summarise}\NormalTok{(}
    \AttributeTok{median\_overlap =} \FunctionTok{median}\NormalTok{(shared),}
    \AttributeTok{.groups =} \StringTok{"drop"}
\NormalTok{  )}

\NormalTok{results}
\end{Highlighting}
\end{Shaded}

\begin{verbatim}
# A tibble: 2 x 2
  pair_type median_overlap
  <chr>              <dbl>
1 ordinary              10
2 wave                  10
\end{verbatim}

\subsection{Step 3: Apply the Decision
Rule}\label{step-3-apply-the-decision-rule}

I set this rule before looking at any real data: \[
\text{Difference} = \text{Median Overlap}_{\text{ordinary}} - \text{Median Overlap}_{\text{wave}}
\]

\begin{Shaded}
\begin{Highlighting}[]
\NormalTok{wave }\OtherTok{\textless{}{-}} 
\NormalTok{  results }\SpecialCharTok{|\textgreater{}}
  \FunctionTok{filter}\NormalTok{(pair\_type }\SpecialCharTok{==} \StringTok{"wave"}\NormalTok{) }\SpecialCharTok{|\textgreater{}}
  \FunctionTok{pull}\NormalTok{(median\_overlap)}

\NormalTok{ordinary }\OtherTok{\textless{}{-}} 
\NormalTok{  results }\SpecialCharTok{|\textgreater{}}
  \FunctionTok{filter}\NormalTok{(pair\_type }\SpecialCharTok{==} \StringTok{"ordinary"}\NormalTok{) }\SpecialCharTok{|\textgreater{}}
  \FunctionTok{pull}\NormalTok{(median\_overlap)}

\FunctionTok{tibble}\NormalTok{(}
  \AttributeTok{wave\_overlap =}\NormalTok{ wave,}
  \AttributeTok{ordinary\_overlap =}\NormalTok{ ordinary,}
  \AttributeTok{difference =}\NormalTok{ ordinary }\SpecialCharTok{{-}}\NormalTok{ wave,}
  \AttributeTok{meets\_rule =}\NormalTok{ ordinary }\SpecialCharTok{{-}}\NormalTok{ wave }\SpecialCharTok{\textgreater{}=} \DecValTok{2}
\NormalTok{)}
\end{Highlighting}
\end{Shaded}

\begin{verbatim}
# A tibble: 1 x 4
  wave_overlap ordinary_overlap difference meets_rule
         <dbl>            <dbl>      <dbl> <lgl>     
1           10               10          0 FALSE     
\end{verbatim}

When meets\_rule is TRUE, the wave pairs kept at least 2 fewer of the
same top 10 pages than the ordinary pairs, which shows the method can
detect the attention shift I planted in the fake data and is ready to be
used on the real protest and Wikipedia data.

If the real data shows the median top 10 overlap in protest wave pairs
is at least 2 pages lower than in ordinary pairs (Difference ≥ 2), I
would conclude that protest waves go along with a shift in what people
read. A smaller difference would mean there is no evidence of a shift,
though even a clear difference would show association, not cause.

\subsection{References}\label{references}

\begin{enumerate}
\def\labelenumi{\arabic{enumi}.}
\tightlist
\item
  Crowd Counting Consortium. \emph{Crowd Counting Consortium Data}.
  Harvard Kennedy School Ash Center and University of Connecticut.
  Harvard Dataverse. \url{https://doi.org/10.7910/DVN/6OPP7H}. About the
  project: \url{https://countingcrowds.org/about/}
\item
  Wikimedia Foundation. \emph{Wikipedia Clickstream}.
  \url{https://dumps.wikimedia.org/other/clickstream/}. Background:
  \url{https://meta.wikimedia.org/wiki/Research:Wikipedia_clickstream}
\item
  Wickham, H., et al.~(2019). Welcome to the tidyverse. \emph{Journal of
  Open Source Software, 4}(43), 1686.
  \url{https://doi.org/10.21105/joss.01686}
\item
  Alexander, R. (2023). \emph{Telling Stories with Data}. Chapman \&
  Hall/CRC. \url{https://tellingstorieswithdata.com/}
\item
  R Core Team. \emph{R: A Language and Environment for Statistical
  Computing}. R Foundation for Statistical Computing.
  \url{https://www.R-project.org/}
\end{enumerate}




\end{document}
